%--------------------------------------------------------------------------------
%
%
%--------------------------------------------------------------------------------
\pagebreak
\section*{LPA — Load Physical Address}
\addcontentsline{toc}{section}{LPA — Load Physical Address}

\begin{iPageDescEntry}{Syntax}
    \texttt{LPA[.D/W/H/B] RegR, ofs ( RegB )} \\
    \texttt{LPA[.D/W/H/B] RegR, RegX( RegB )}
\end{iPageDescEntry}

\begin{iPageDescEntry}{Format}
The \texttt{LD} instruction uses the indexed and register-indexed instruction formats. \\

    \begin{flushleft}
    \drawInstrFormat
        {0,1,3,5,6.5,8.5,9.5,16}
        {\OpCodeGrpSYS, \OpCodeLPA, RegR, 0, RegB, dw, ofs}
    \end{flushleft}

    \begin{flushleft}
    \drawInstrFormat
        {0,1,3,5,6.5,8.5,9.5,11.5,16}
        {\OpCodeGrpSYS, \OpCodeLPA, RegR, 1, RegB, dw, RegX, 0}
    \end{flushleft}
\end{iPageDescEntry}

\begin{iPageDescEntry}{Description}
    The \texttt{LPA} (Load Physical Address) instruction computes the physical 
    address corresponding to the given virtual address and loads it into 
    \texttt{RegR}. If there is no translation, a value of zero is returned. Two formats are available: the immediate-offset format and the register-indexed format. 
    
    In the immediate-offset format, the signed immediate offset is shifted right by the \texttt{dw} field of the instruction and added to \texttt{RegB} to form the effective address.  
    In the register-indexed format, \texttt{RegX} is added to \texttt{RegB} to form the effective address.

    The physical address, if found, is returned from the information in the TLB. If there is a separate instruction and data TLB is used. No memory access is performed. For direct mapped segments the virtual address is the physical address.
    
    Optionally, the computed effective address may also be written back to the base register \texttt{RegB}. When this option is enabled, the sign of the offset determines the update behavior:  
    – A negative offset causes the base register itself to be used as the fetch address.  
    – A positive offset causes the sum of the base register and offset to be used.  
    In both cases, the base register is updated.
\end{iPageDescEntry}

\begin{iPageDescEntry}{Operation}
	\texttt{RegR <- tanslateAdr( memOperand( instr )));} 
\end{iPageDescEntry}

\begin{iPageDescEntry}{Exceptions}
	- Alignment Trap \\
    - Privileged instruction trap. \\
    - Non access data TLB trap.
\end{iPageDescEntry}

\begin{iPageDescEntry}{Notes}
    Useful for operating systems to probe or manage virtual to physical mappings. \\
 \end{iPageDescEntry}
